\documentclass[output=paper]{langscibook}
\ChapterDOI{10.5281/zenodo.18845934}
\author{Marilyn M. Vihman\orcid{}\affiliation{University of York (Language and Linguistic Science), UK; University of California, Berkeley (Linguistics), US} and Virve-Anneli Vihman\orcid{}\affiliation{University of Tartu, Estonia}}
\title[Linguistic interaction in early bilingual development]{Linguistic interaction in early bilingual development: A review of studies examining single- and mixed-language utterances}
\abstract{Over the past 30 years or so a steady stream of papers has addressed bilingual development and the possible effects – or lack of effects – of cross-linguistically differing structures on a bilingual child’s acquisition of those structures. Drawing generalizations has proven difficult, partly due to the many variables that enter in. Furthermore, the theoretical assumptions of the authors vary from one study to the next; this can make it difficult to evaluate the different conclusions drawn. A few studies have specifically focused on the long-standing puzzle of how best to weigh the effects of (i) experience with adult input patterns and (ii) the child’s own relative proficiency in each language on the diverse manifestations of bilingual child usage, including code-switching. Here we provide an overview of the findings of the bilingual child literature on the extent of linguistic interaction in the acquisition of morphosyntactic structure. We review 47 studies of the spontaneous language use of bilingual children up to age 4 and emphasise the point that child production is a moving target that necessarily reflects the dynamics of the child’s own emergent grammatical systems and possible shifts in their exposure to and use of the languages in question.}
\IfFileExists{../localcommands.tex}{
  \addbibresource{../localbibliography.bib}
  \usepackage{tabularx,multicol}
\usepackage{url}
\urlstyle{same}

 
\usepackage{langsci-optional}
\usepackage{langsci-lgr}
\usepackage{langsci-gb4e}

\usepackage{multirow}
\usepackage{cleveref}
\usepackage{subfigure}
\usepackage{langsci-branding}

  \input{../localcommands}
  \input{../localhyphenation}
  \togglepaper[6]%%chapternumber
}{}

\begin{document}
\maketitle
%\shorttitlerunninghead{}%%use this for an abridged title in the page headers


\section{Introduction}

A central concern of studies of bilingual language development is the nature and extent of interaction – also termed interference (e.g., \citealt{ParadisNavarro2003}), transfer (e.g., \citealt{YipMatthews2000}) or influence (e.g., \citealt{Serratrice2013}) – between the linguistic structures of the two languages to which the child has been exposed. This is crucial for understanding the process of bilingual acquisition, in which the child’s neural, cognitive and linguistic systems develop in interaction with two grammatical systems. Clarifying how cross-linguistic interaction operates in bilingual children is also important for better understanding adult bilingual processing, and for adapting models of bilingual language use.

Evidence of interaction can most readily be drawn from the child’s speech production; such interaction indicates an effect of bilingual experience on the child's developing grammatical systems. Here we provide an overview of studies from the last three decades that address cross-linguistic interaction in bilingual children’s production during the period when morphosyntax emerges. We understand code-switching (or the production of ‘mixed-language utterances’) to constitute one manifestation of such interaction. Note, moreover, that the evidence we investigate from mixed utterances is not the inclusion of elements from more than one language \textit{per se}, but rather the effect of the structure of one language on production in the other, the same kind of evidence we investigate in single-language utterances. As cross-linguistic influence may manifest differently in utterances with lexemes from one language compared to mixed-language utterances, we divide our sample of studies into (i) those that set mixed utterances aside to focus on cross-linguistic influence in single-language utterances and (ii) those that specifically address utterances that include both languages.

Many factors conspire to make the effects of linguistic interaction on bilingual language use a challenging research topic. In adults such factors include the difficulty of measuring and controlling for heterogeneity in proficiency, usage and relative dominance in the two languages, differences in their typological relatedness and also task differences that might affect executive control, cross-language cues and language activation (see \citealt{Grosjean1998,BaumTitone2014}). For developmental bilingualism, a critical issue is the \textit{extent of exposure,} or \textit{opportunities for hearing and using each language} – and this is liable to \textit{change over time} (cf. \citealt{TullochHoff2022}). \textit{Language dominance} – the child’s relative proficiency in each language – is affected by relative exposure to and opportunities for use of the two languages but also interacts with the child’s \textit{overall cognitive and linguistic development}. In young children, furthermore, it is important to consider the wide range of \textit{individual differences,} which are also commonly observed in monolingual development (cf. \citealt{Serratrice2005}, who explores the effect of individual differences on learning strategy in two bilingual children raised in similar circumstances, though with different language pairs).

Beyond these sources of complexity for analysis, the differing \textit{theoretical assumptions} of authors also make it difficult to draw generalizable conclusions from disparate studies. We note further that \textit{language usage in the community} is seldom taken into account, with the exception of a few recent studies (e.g., \citealt{YipMatthews2016}: Cantonese \& English in Hong Kong; \citealt{QidiBiase2020}: Mandarin and English in Australia; \citealt{BalamEtAl2021} and \citealt{vanOschEtAl2023}: Spanish and English in Miami, Florida and the Netherlands, resp.; \citealt{PhillipsDeuchar2022}: Welsh and English in North Wales).  The linguistic context can be taken to include family language policy, community norms, the extent of bilingual usage in the community, community attitudes toward bilingualism and/or toward each language and the pattern of interaction of the two languages in the community. For example, \citet{YipMatthews2016} mention that although code-switching is widespread in Hong Kong, families there vary in extent of switching, and this influences the children’s usage. \citet{PhillipsDeuchar2022} specifically investigate how the child’s pattern of switching is related to the prevailing adult usage in a bilingual Welsh-English community, which resembles Hong Kong in the prevalence of switching. The linguistic context, in both the home and the larger community, is relevant for the question of cross-language influence as well as for understanding developmental code-switching patterns.

\section{Background}

Over the past 30-odd years the number of investigators interested in documenting and discussing the process of learning two languages in early childhood has expanded considerably, and the study of bilingual language development has surged. Our initial search on Google Scholar (February 2023), using the simple term ‘bilingual language development’ and limited to studies from 1990 on, yielded nearly 5000 titles. Kroll \& Bialystok (2013, p. 498) document a ``steep and continuing rise in both publications and citations" on bilingualism in general from 2003 on; we see a similar increase in bilingual development, with half of the responses from our initial search dating to 2010 or later. To address the question of cross-linguistic interaction in bilingual children in a systematic, quantitative manner, \citet{VanDijkEtAl2022} recently provided a meta-analysis of experimental studies, restricted to those with explicit measurement and at least two bilingual and two monolingual children. All participants were aged three years or older. That paper examines evidence for cross-linguistic influence from 26 studies, with data from 750 children and 17 unique language pairs (Hebrew is the only non-European language included). These studies typically involve elicited production, with a few testing grammaticality judgments and a small number involving comprehension only. Four-year-olds were the group most often tested, but school-age children up to the age of ten were also included.

The \citeauthor{VanDijkEtAl2022} study is systematic and thorough, but the focus on older children and experimental results leaves a significant gap. Younger children may be more variable and difficult to include in well-controlled experiments but assessing the extent of separation or cross-linguistic interaction in their emergent grammars is essential for forming an adequate picture of the trajectory of bilingual grammatical development. For very young children, spontaneous production provides the most reliable information, reflecting their behavior in a natural, everyday context. Here, then, we provide a complementary overview of studies investigating data (i) from the early years of morphosyntactic development, (ii) in spontaneous speech production. Considering children’s early use of their developing languages should enable us to detect the effect of the relative accessibility of structures, their parallel development and the extent of bilingual interaction as the two grammatical systems first emerge and gradually come into regular use. To achieve these ends we focus on studies that specify cross-linguistic effects on particular morphosyntactic structures, in either single- or mixed-language utterances. 

In a second search, following \citet{VanDijkEtAl2022}, we entered ‘bilingual child’ as our basic search term, combined with \{‘cross-language’\} or \{‘cross-linguistic’\} and \{‘effects’, ‘transfer’, ‘influence’, ‘interaction’, ‘activation’\}. Together these search terms yielded over 9000 titles, with a certain amount of overlap. We then refined the results by restricting our sample to studies providing \textit{data from spontaneous speech} (based on transcripts of audio- or video-recordings, diary notes, or both) and relating to \textit{morphosyntax}. In addition, for comparability, we chose sources that provide data derived from (i) \textit{typically developing children,} (ii) recorded from \textit{the} \textit{age of four or younger,} (iii) exposed to \textit{two oral languages} and (iv) including \textit{morphosyntactic analysis of multiword utterances}, whether those include code-switching or not. We thus excluded studies using elicited data or focusing solely or mainly on phonetics, phonology, the lexicon, sign languages or bimodal bilingualism, as well as studies deriving data from children with atypical development; we also excluded overviews, whether purely theoretical or oriented toward teaching or clinical use. Finally, we included only studies that provide \textit{explicit analysis and discussion of morphosyntactic interaction}, in single- or mixed-language utterances, and with specific reference to the morphosyntactic structures affected. With those restrictions, and with the addition of further references derived from the readings, we identified 47 empirical studies. Note, further, that our search was restricted to articles published in English, admittedly a limitation for a study seeking information on typologically varied language pairs.

Linguistically, the studies are wide-ranging: Overall, 21 different languages and 30 unique language pairs are included in at least one study each. English is one of the languages of a pair in 34 studies, German in 23, French or Italian in 15 each, Spanish in 10; note that some studies report and analyze data from more than one language pair and also that some of the same children’s data are treated in more than one study, with a focus on different structures. The 47 studies include 197 children altogether, or a mean of 4.2 children per study. There are 30 single-child case studies; 14 studies include at least four children learning one or more of the language pairs they address; only three include more than ten. In short, these studies, selected for their analyses of early spontaneous bilingual language use, generally report data from very few children.\footnote{\citet{SağınŞimşekAntonovaÜnlü2021} is the only study included here to have used elicitation – namely, of the Frog Story narrative \citep{Mayer1969} – as part of their (annual) data collection; this was supplemented with recordings of naturalistic play. The children – four monolinguals and four in each of three bilingual groups (English-, German- and Russian-Turkish) – ranged in age from 2;11 to 3;11 in the first recordings, at the upper limit of our sample.} The review is largely qualitative, given both the type of data available for younger children and the limited numbers of participants per study.

At least three key issues in the literature on bilingual language development have sharply changed over the last three decades. First, attitudes and assumptions have changed:  {In 1988} Petersen still felt it necessary to insist that her daughter’s addition of English inflections to Danish content nouns did not reflect \textit{‘random lexical choices’} or \textit{‘confusion’} (as the popular view might hold: \citealt{Goodz1989}), but instead could be seen to reveal \textit{systematic patterning}. Denials that bilingual children are ‘confused’ have continued to surface from time to time (e.g., \citealt{YipMatthews2000}: 207; \citealt{ParadisNavarro2003}: 372), but there has long been general agreement on the systematicity of child code-switching, in accord with the consensus, securely established over this same period, that adult code-switching is itself systematic (cf., e.g., \citealt{Kroll2005,Deuchar2012,Deuchar2020,ValdesKroffDussias2023}).

Secondly, in the late 1980s the question arose as to whether the bilingual child \textit{learns from the input} to use intraclausal code-switching, or code-mixing (the ‘modeling hypothesis’: \citealt{ComeauEtAl2003}). \citet{Goodz1989} reported that, in a longitudinal study of four families in Montreal, the input the child heard tended to include switching, despite the parents’ stated commitment to each of them using their native language only. The proportion of mixed language utterances was low, however, and tended to be initiated by the child and then echoed by the parent rather than the reverse. In a study designed to test the question, \citet{GeneseeEtAl1995} found no evidence that the language mixing of the five bilingual children they recorded with each parent was correlated with the extent of mixing in their parents’ speech. Mixing in parental input now seldom seems to be invoked as a source of child code-switching (see \citealt{MengMiyamoto2012,TullochHoff2022}).

A closely related aspect of the issue of learning from the input is the question of parental strategy. Over half of the studies we review here (27 of 46) report that the participating families observe the \textit{‘one person-one language’} principle in raising a bilingual child (as did the families recorded in the studies of both Goodz and Genesee et al. in Montreal). The principle was originally recommended by \citet{Ronjat1913}; it was later forcefully promoted by investigators such as \citet{Dopke1992} and \citet{Dehouwer1995,Dehouwer2005}. Yet in the papers we review parental strategy is seldom mentioned as a factor contributing to either linguistic interaction or the lack of it.  One exception is \citet{YipMatthews2016}, who report a higher level of mixing in the production of children from families freely using both Cantonese and English in the home, with considerable adult mixing, as compared to those raised in homes following the ‘one person, one language’ principle, which is less common in present-day Hong Kong. It is worth noting that where the home language differs from that of a uniformly majority-language community, many families consider primary use of the minority language in the home to be essential to ensure bilingual development (e.g., \citealt{DeucharQuay2000,Vihman1985,Vihman1998}; cf. also \citealt{Gawlitzek-MaiwaldTracy1996}: 908).

In general, the pattern of \textit{language use in the community} should be taken into account (\citealt{PhillipsDeuchar2021}). \citet{Byers-Heinlein2013}, a large-scale study of parental input and its effect on vocabulary growth, reports considerable mixing in bilingual families in Vancouver, most of whom are part of a large bilingual community. Beyond that, \citet{PhillipsDeuchar2021} show that in North Wales, where the minority language (Welsh) consistently provides the matrix into which content words from the majority language (English) are inserted in the commonly occurring code-switching of adults, the children’s mixed utterances follow suit; the reverse pattern – with English as the language of the verb and function words, Welsh as the source of content-word insertions – is virtually unattested in the corpus. Attention to the language of the community might also help to account for \citegen{Petersen1988} findings (cf. also \citealt{Lanza1997,Lanza_language_1997-1,Vihman1998}, for example). Where the language of the home differs from that of the community, the ‘matrix language’ – the language in which the predicate and most of the function words are expressed (\citealt{Myers-Scotton1993duelling}) – is most likely to be that of the home, while the community language serves as a source of content-word insertions, at least in the preschool years.

Finally, the assertion that children are \textit{able to differentiate between their two grammars} from the beginnings of word combination, if not sooner, runs as a thread through much of the literature. \citet{Genesee1989} appealed to the considerable evidence of infants’ precocious capacity for distinguishing between speech sounds to shed doubt on the claim that bilingual children go through a stage of treating their two languages as if they constituted a single unified or ‘fused’ system (\citealt{Volterra1978}; see \citealt{Vihman2014}, Ch. 3, for a review of the extensive infant speech perception literature supporting Genesee’s point, and \citealt{Byers-Heinlein2014}, for a reassessment of the issue of early language differentiation). More recent perception studies have provided ample experimental evidence that infants are able to discriminate  between languages even at birth (see \citealt{Nazzi1998,Ramus2002,Byers-HeinleinEtAl2010}), and that bilingual children can segment speech in both their languages by the middle of the first year, given sufficient exposure (e.g., \citealt{BoschEtAl2013,Polka2017}). Accordingly, children must be sensitive from their earliest months to the presence of more than one language in their environment.

The claim that the child has ‘two systems from the start’ is also rooted in the theoretical clash that has pervaded acquisition studies since the publication of Chomsky’s \textit{Aspects of the Theory of Syntax} (1965), with its proposal of an innate Language Acquisition Device. If, as generative approaches have traditionally maintained, children set the parameters appropriate for each language by accessing Universal Grammar, then the two linguistic systems should develop autonomously; according to this view, there is no reason to expect interaction, at least not at the level of ‘competence’, or the child’s underlying knowledge (see \citealt{ParadisGenesee1996}, who defend this position).

The claim of independence for each of a child’s grammatical systems continues to be reiterated in the literature, even though no one has defended the ‘fused system’ hypothesis for over 30 years. As an example of current views, \citet{Kupisch2007} provides evidence of interaction in the acquisition of Italian and German determiners that goes against the assumption that bilingual children have ‘two autonomous grammatical systems’, but she maintains that language influence (or interaction) and language separation are not mutually exclusive (p. 76). Most investigators today would probably agree with this more nuanced conclusion. 

An alternative perspective on the issue, from a usage-based theoretical position, is provided by  \citet{GaskinsEtAl2021}. Each of a child’s experiences with hearing and producing language leaves its memory trace, or set of exemplars, in the child’s mind and thus adds to the child’s stock of linguistic knowledge. Previously experienced lexical or phrasal units provide potential material for linguistic expression. By this account, both frequency-based entrenchment and ‘priming’, or the experience of the immediately preceding speech, may explain a child’s recourse to units from one language or the other or from both within the same utterance. \citet{GaskinsEtAl2022} offer some suggestions for future research to test the validity of these ideas. In effect, most researchers might now agree that the systems emerge and develop in parallel in a child with sufficient exposure to and opportunities for use of two languages; thus some degree of linguistic interaction can be expected.

One way to understand how linguistic systems may interact in the child is the suggestion that the acquisition of a more easily learned construction in one language supports learning of the equivalent structure in the other. The proposal of linguistic facilitation, which \citet{BernardiniSchlyter2004} termed the ‘ivy hypothesis’, derives from evidence that bilingual children sometimes benefit from \textit{accelerated} \textit{learning} of aspects of grammar that are comparable in their two languages (e.g., \citealt{Kupisch2007}). In contrast, the identification of difficulty with a construction as compared (by proportion of recorded attempts) with monolingual children can suggest some \textit{delay} due to the child’s experience of two systems. The errors in monolingual children’s production can provide a useful developmental comparison point for bilingual children’s production. However, in cases where this involves comparing a single bilingual and a single monolingual child, the findings are not generalizable, given the inherent variability in age and trajectory of individual children’s linguistic advance.

\citet{RothmanEtAl2023} problematize the use of monolingual children as controls in studies of bilingual development, observing that, for example, a monolingual ‘homeland’ speaker group is not necessarily a revealing comparison for heritage language speakers. They note, however, that ``if one is interested in documenting the (potential) role that crosslinguistic influence has in the development of bilingual grammars in childhood, it could be reasonable to compare a child bilingual group to a monolingual one" (2023: 319), as long as it is motivated by the research question. The small sample size in the studies discussed here means that any comparisons are indicative rather than conclusive. Furthermore, where a comparison indicates delay, this may derive as much from limited input or exposure as from cross-linguistic interaction. For identifying interaction in small groups of participants, it might be more fruitful to compare bilingual groups with differing language pairs.

Another type of evidence of interaction derives from a bilingual child’s mixed-language utterances, which often show reliance on function words or inflections from one language along with content words from the other; such usage suggests memory or representational challenges. Note that the mixing of content words from the larger community need not mean that the child has not yet learned the relevant words in the home language or, in other words, that there is an actual ‘lexical gap’; instead, such mixing may result from a momentary ‘retrieval problem’ (\citealt{Lanvers2001}: 449) or ‘retrieval failure’ (\citealt{Serratrice2005}: 169). Certainly momentary lapses in access are not unknown in adult speakers, whether bilingual or not. Such lapses likely account for much of the variability in child production at every level.

We focus here on morphosyntactic interaction in early bilingual child production, dividing our overview into the two major strands we find in the literature: (i) evidence of interaction in single-language utterances in either language and (ii) evidence from code-switching or ‘mixed-language’ utterances. Our research questions are: 

\begin{enumerate}
\item \textit{What aspects of the grammar} are most likely to be affected, or \textit{which linguistic structures} are involved in cross-linguistic interaction in the young bilingual child?
\item What \textit{types of interaction} have been reported? And to what extent does \textit{dominance} (or greater proficiency) in one language account for the effects?
\item What role is played by the \textit{relative accessibility and compatibility of structures?}
\end{enumerate}

\section{Morphosyntactic interaction between linguistic systems in early bilingual child production} 

Our division of the studies into those addressing morphosyntactic effects in single-language utterances and those addressing intraclausal mixing (including intralexical mixing, between the stem and a bound morpheme) will give an idea of the ways in which interaction is manifested.

\subsection{Evidence from single-language utterances} \label{sec:vihmannvihmann:3.1}

Thirty-eight (81\%) of the 47 studies in our sample address the issue of morphosyntactic interaction in single-language utterances. Since the number of children in most of these studies is small, we group the studies together by the linguistic structures they analyze, to the extent possible (Tables 1–8, according to topics investigated); some studies appear under more than one topic (and some studies also consider mixed-language utterances). The tables indicate the linguistic structures, the language pairs included, the language of the children’s community (with the country or region where data were collected, if specified), and the number of children per language pair. Each study is marked for the type of interaction reported, if any: \textit{acceleration}, \textit{delay} – \textit{not} including delay said to be due to \textit{lack of sufficient experience} with the language in question – or \textit{transfer} (effects other than acceleration or delay). The studies are listed chronologically within each linguistic structure. 

Under Dominance, we indicate whether the authors find it to be a factor in the linguistic interaction they identify; N/A indicates that dominance does not figure in the analyses. Relative ‘dominance’ or ‘balance’ was established in different ways in the 20 studies that mention it. Eleven studies base their characterization on MLU counts. \citet{Kupisch2007} includes additional measures (type/token ration, upper bound, or longest utterance, number of noun and verb types, number of utterances in 30 minutes). Others also consider aspects of the child’s input in each language, based on attendance at school (\citealt{Austin2009}) or nursery (\citealt{Blom2010}), parental questionnaire (\citealt{FernándezFuertesLiceras2010,LicerasEtAl2012}) or the child’s daily activities with different interlocutors (\citealt{QidiBiase2020}). Five more studies calculate MLU but do not base their judgment of balance on it. For example, \citet{Dopke1998,Döpke2000} and \citet{HulkMüller2000} make use of MLU to establish developmental phases rather than dominance. \citet{Mishina-Mori2020} calculates MLU but uses number of word types to establish dominance. Finally, both \citet{AdnyaniEtAl2018} and \citet{Silva-corvalan2008} are diary studies that rely on informal observation of the child’s experience of the two languages to assess relative dominance. In short, language use by both the children and their interlocutors is variously considered in these studies, with varying degrees of rigor.

Despite differences in methods, analyses and interpretations, 25 of the 38 studies listed in Tables 1–8 (66\%) find evidence of interaction in at least one of the structures they analyse.\footnote{Note that although we list studies repeatedly in the tables where more than one language pair was studied, we count the studies only once each in regards to their findings on interaction and tne role of dominance.} The studies that address \textit{copula selection or omission} (\tabref{tab:vihmannvihmann:1}) and \textit{object omission} (\tabref{tab:vihmannvihmann:2}) all provide evidence of interaction.


\begin{table}
\begin{tabularx}{\textwidth}{QQ@{}rcQ}
\lsptoprule
{ Study} & { Languages}   (location)& { N}  & { Dominance} & { Type of Interaction}\\
\midrule
\citet{Silva-corvalan2008} & \textbf{English} Spanish (United States) & 1 & + & Delay in acquisition of Spanish \textit{ser/estar} \\
\tablevspace
\citet{FernándezFuertesLiceras2010,LicerasEtAl2012} & English \textbf{Spanish} (Spain) & 2 & - & Acceleration:  Less copula omission in English (due to more complex Spanish system)\\
\lspbottomrule
\end{tabularx}
\caption{Studies investigating cross-linguistic interaction in \textbf{copula} \textbf{selection} \textbf{or} \textbf{omission} in single-language utterances. The community language(s) are marked in bold. We indicate whether the authors see Dominance as a factor affecting the interaction (+ / -); NA indicates that the study does not mention dominance.}
\label{tab:vihmannvihmann:1}
\end{table}


\begin{table}
\begin{tabularx}{\textwidth}{QQrcQ}
\lsptoprule
{ Study} & { Languages}  { (location)} & { N}  & { Dominance} & { Type of Interaction}\\
\midrule
\citet{HulkMüller2000}, \citet{Muller2001} & \textbf{Dutch} French \mbox{(The Netherlands)} & 1 & - & Delay in object provision in Romance languages \\
%\hhline%%replace by cmidrule{--~--}
& French \textbf{German} (Germany) & 1 &  & \\
%\hhline%%replace by cmidrule{~-~~~}
& \textbf{German} Italian (Germany) & 2 &  & \\
\tablevspace
\citet{YipMatthews2000} & \textbf{Cantonese} \textbf{English} (Hong Kong) & 1 & + & Delay in object provision in English\\
\tablevspace
\citet{Mishina-Mori2020} & \textbf{English} Japanes (United States) & 2 & - & Delay in object provision in English\\
\lspbottomrule
\end{tabularx}
\caption{Studies investigating cross-linguistic interaction in \textbf{object} \textbf{omission} in single-language utterances.}
\label{tab:vihmannvihmann:2}
\end{table}

In contrast, of the six studies looking at \textit{inflection or subject-verb agreement} (\tabref{tab:vihmannvihmann:3}), only two report interaction.

\begin{table}[t]
\begin{tabularx}{\textwidth}{QQ@{}rcQ}
\lsptoprule
{ Study} & { Languages} { (location)} & { N}  & { Dominance} & { Type of Interaction}\\
\midrule
\citet{ParadisGenesee1996} & English \textbf{French} (Montreal, Canada) & 3 & NA & None\\
\tablevspace
\citet{SinkaSchelleter1998} & Latvian \textbf{English} (England) & 1 & NA & None\\
\tablevspace
\citet{Austin2007} & \textbf{Basque} \textbf{Spanish} (Spanish Basque country) & 20 & NA & Delay in Basque ergative use  \\
\tablevspace
\citet{HacohenSchaeffer2007} & English \textbf{Hebrew} (Israel) & 1 & - & None\\
\tablevspace
\citet{Soriente2014} & \textbf{Indonesian} Italian (Indonesia) & 1 & + & None\\
\tablevspace
\citet{AdnyaniEtAl2018} & German \textbf{Indonesian} (Indonesia) & 1 & + & Transfer: Vocative + infinitive predicates \textit{pro} subject + finite verb in German\\
\lspbottomrule
\end{tabularx}
\caption{Studies investigating cross-linguistic interaction in \textbf{inflection,} \textbf{or} \textbf{subject-verb} \textbf{agreement} in single-language utterances}
\label{tab:vihmannvihmann:3}
\end{table}

Of the remaining studies focusing on inflection, \citet{SinkaScheletter1998}, in a case study of Latvian-English acquisition, find no cross-linguistic effects, but provide neither reference to baseline nor monolingual data. Emphasizing the need for evidence of systematic differences to support a claim of interaction, \citet{ParadisGenesee1996} review studies of monolingual acquisition in French and English that demonstrate relatively earlier mastery of the French system of inflectional marking, due in part to the more central role of inflection in French. Their subsequent analysis of the acquisition of finiteness and agreement in three French-English bilingual children aged 2-3 years shows that they acquire the two systems in much the same way as monolinguals, with the children expressing finite verbs in French before they begin to do so in English (although the bilinguals do produce more finite utterances in English overall than is typical of monolinguals). \citet{HacohenSchaeffer2007} compare the Hebrew production of a single Hebrew-English-learning child with five Hebrew monolingual learners. They find an equally low percentage of subject-verb agreement errors in all the children. And \citet{Soriente2014} finds a delay in the acquisition of Italian inflection but considers it the result of insufficient input in that language.

As regards \textit{wh-movement} (\tabref{tab:vihmannvihmann:4}), the two studies that find no interaction – \citet{Park-Johnson2017} and \citet{QidiBiase2020} – both concern early sequential bilingualism.

\begin{table}
\begin{tabularx}{\textwidth}{QQ@{}rcQ}\lsptoprule
{ Study} & { Languages} { (location)} & { N}  & { Dominance} & { Type of Interaction}\\
\midrule
\citet{YipMatthews2000} & \textbf{Cantonese} \textbf{English} (Hong Kong) & 1 & + & Delay in English WH-fronting \\
\tablevspace
\citet{Mishina-mori2005} & \textbf{English} Japanese (United States) & 2 & NA & Delay: Overuse of WH-fronting in Japanese\\
\tablevspace
\citet{Soriente2007} & Italian \textbf{Indonesian} (Indonesia) & 1 & + & Delay in acquisition of WH-order in Italian\\
\tablevspace
\citet{Park-Johnson2017} & \textbf{English} Korean (United States) & 3 & NA & None\\
\tablevspace
\citet{QidiBiase2020} & \textbf{English} Mandarin (Australia) & 1 & - & None\\
\tablevspace
\citet{YipMatthews2000} & \textbf{Cantonese} \textbf{English} (Hong Kong) & 1 & + & Delay in English WH-fronting \\
\lspbottomrule
\end{tabularx}
\caption{Studies investigating cross-linguistic interaction in \textbf{WH-movement} in single-language utterances}
\label{tab:vihmannvihmann:4}
\end{table}

The children observed in \citeauthor{Park-Johnson2017}’s study initially experienced mainly Korean input, not only in the home but also in church and weekend Korean language school. The children began to learn and use English, the language of the larger (American) society, at nursery school. In the case of \citeauthor{QidiBiase2020}, the child’s home (in Australia) included monolingual Mandarin-speaking grandparents in addition to the parents and an aunt. The child did not actively begin \textit{producing} English until after six months of silent attendance at a childcare center, beginning at 2;8. In these studies, then, the children were mainly exposed to the home language up to about age 3, when they began to learn and use English in a nursery setting. Such early successive bilingualism may be qualitatively different from bilingual exposure and use from the start (\citealt{Vihman1982}).

Eight studies addressed \textit{subject provision} (\tabref{tab:vihmannvihmann:5}) in languages that do not require it (and in contexts in which overt pronominal subject provision is pragmatically marked): Six of these found evidence of overprovision of subjects in the children’s pro-drop language, including one study of Spanish-English acquisition (\citealt{ParadisNavarro2003}, who deemed their findings inconclusive), while two other studies of Spanish-English acquisition (\citealt{LicerasEtAl2012,Villa-GarciaSuarez2016}) did not.\footnote{Three of the four bilinguals included in \citet{Villa-GarcíaSuarez2016} were also included in either \citet{ParadisNavarro2003}, or \citet{LicerasEtAl2012}.} \citet{Blom2010} finds underprovision of obligatory subjects in Dutch, an effect of Turkish, but only in the case of children for whom Dutch is the weaker language.

\begin{table}\small
\begin{tabularx}{\textwidth}{QQ@{}rcQ}
\lsptoprule
{ Study} & { Languages}  { (location)} & { N}  & {{ Dominance}} & { Type of Interaction}\\
\midrule
\citet{ParadisNavarro2003} & \textbf{English} Spanish (England) & 1 & {NA} & Delay: Overprovision of subjects in Spanish\\
\tablevspace
\citet{Serratrice2004} & \textbf{English} Italian (Scotland) & 1 & {NA} & Delay: Overprovision of pronominal subjects in Italian\\
\tablevspace
\citet{HacohenSchaeffer2007} & English \textbf{Hebrew} (Israel) & 1 & {-}
& Delay: Overprovision of subjects in Hebrew\\
\tablevspace
\citet{Blom2010} & \textbf{Dutch} Turkish (The Netherlands) & 4 & {+} & Delay: Under-provision of subjects in Dutch (when Turkish-dominant)\\
\tablevspace
\citet{Haznedar2010} & English \textbf{Turkish} (Turkey) & 1 & {NA} & Delay: Overprovision of subjects in Turkish \\
\tablevspace
\citet{LicerasEtAl2012} & English \textbf{Spanish} (Spain) & 2 & {-} & None\\
\tablevspace
\citet{SchmitzEtAl2012} & \textbf{German} Italian (Germany) & 3 & NA & {Delay: Overprovision of pronominal subjects in German}\\
                       & French \textbf{German} (Germany) & 4 &  &   \\
%\hhline%%replace by cmidrule{~--~~~}
                      & \textbf{French} Italian (France) & 1 &  &   \\
\tablevspace
\citet{Villa-García2SuárezPalma016} & \textbf{English} Spanish (USA, England) & 2 & NA & {None}\\
                              & English-\textbf{Spanish} (Spain) & 2 &  &   \\
%\hhline%%replace by cmidrule{-----~}
\lspbottomrule
\end{tabularx}
\caption{Studies investigating cross-linguistic interaction in \textbf{Subject} \textbf{provision} in single-language utterances}
\label{tab:vihmannvihmann:5}
\end{table}

Five of the eight papers investigating aspects of \textit{word order} report evidence of cross-linguistic effects, while three do not (\tabref{tab:vihmannvihmann:6}).

\begin{table}
\begin{tabularx}{\textwidth}{>{\raggedright}p{35mm}@{~~}Q@{}rcQ}
\lsptoprule
{ Study} & { Languages} { (location)} & { N}  & { Dominance} & { Type of Interaction}\\
\midrule
\citet{ParadisGenesee1996} & English \textbf{French} (Montreal, Canada) & 3 & NA & None\\
% \tablevspace
\citet{Dopke1998} & { \textbf{English} German} (Australia) & 3 & - & Delay in acquisition of word order in German VP \\
% \tablevspace
\citet{SinkaSchelleter1998} & { \textbf{English} German} (England) & 1 & NA & None\\
% \tablevspace
\citet{Döpke2000} & { \textbf{English} German} (Australia) & 4 & - & Delay in acquisition of word order in German VP \\
% \tablevspace
\citet{Hauser-GrüdlEtAl2010} & \textbf{German} Italian (Germany) & 1 & NA & Acceleration: Earlier adherence to German main-clause V2 order\\
% \tablevspace
\citet{AdnyaniEtAl2018} & German \textbf{Indonesian} (Indonesia) & 1 & + & Transfer: Use of German VP order in Indonesian \\
% \tablevspace
\citet{AndersenBentzen2018} & { English \textbf{Norwegian}} (Norway) & 3 & - & Transfer: Overuse of V2 in English \\
% \tablevspace
\citet{SağınŞimşekAntonovaÜnlü2021} & { \textbf{English} Turkish} (England) & 4 & NA & None\\
%\hhline%%replace by cmidrule{--~--}
& { \textbf{German} Turkish} (Germany) & 4 &  & \\
%\hhline%%replace by cmidrule{~-~~~}
& { \textbf{Russian} Turkish} (Russia) & 4 &  & \\
\lspbottomrule
\end{tabularx}
\caption{Studies investigating interaction in \textbf{word} \textbf{order} in single-language utterances}
\label{tab:vihmannvihmann:6}
\end{table}

\citet{ParadisGenesee1996} compare the bilingual children’s production in the two languages as regards placement of the negator, which follows the finite verb in French but precedes it in English. Although the data from their three participants is sparse, these investigators do find some evidence of English influence on the children’s French, such as pre-verbal placement of the negator or the use of \textit{non} in lieu of \textit{pas} (an apparent calque on English \textit{not)}. However, they dismiss ‘these aberrant examples’ as ‘most likely performance errors’ (p. 15). 

\citet{SinkaSchelleter1998} find no evidence of cross-linguistic effects on noun-verb placement in a child learning English and German, in the UK, at 2;0-2;6, while \citet{Dopke1998,Döpke2000} reports a delay in the acquisition of German VP order under the influence of English in four children learning English and German in Australia. Analyzing data from earlier studies of one bilingual Italian-German and one monolingual German child, \citet{Hauser-GrüdlEtAl2010} find acceleration in German under Italian influence: The monolingual German learner they observe goes through a phase of using, in main clauses, the verb-final word order pattern found in German subordinate clauses; this persisted to about age 2;8. Thereafter, the child gradually uses more verb-second order (V2) in main clauses. In contrast, the bilingual child fails to show early use of V-final order, producing V2 from early on instead. ``Since in Italian, V-final word order patterns hardly exist, we take this lack of V-final in Lukas to be due to cross-linguistic influence of his Italian L1 onto his German L1." (p. 2643)

The four children in each of the three bilingual groups in \citet{SağınŞimşekAntonovaÜnlü2021} showed, in their Turkish production, a developmental trajectory distinct from that of the four Turkish monolinguals. Specifically, all the children make use of the canonical Turkish SOV order in the earliest recordings, but the monolinguals make increasing use of pragmatically based non-canonical orders as they grow older (up to 30\% by age 5;5), whereas the heritage-language Turkish learners continue to make almost exclusive use of SOV. The authors argue that a lack of sufficient exposure to Turkish input can best account for the difference; they find direct cross-linguistic influence to be implausible, given that the Russian bilingual learners, who are exposed to more variable word orders, show no better ability to adopt the kind of pragmatic variation Turkish allows than do the children acquiring English or German, with their more fixed verb/object orders.

\begin{table}[b]
\begin{tabularx}{\textwidth}{QQ@{}rcQ}
\lsptoprule
{ Study} & { Languages} { (location)} & { N}  & { Dominance} & {{ Type of Interaction}}\\
\midrule
% &  & 1 & & \\
\citet{HulkMüller2000} & \textbf{Dutch} French (Netherlands) & 1 & – &{None}\\
%\hhline%%replace by cmidrule{~--~~~}
& \textbf{German}~Italian (Germany) & 1 &  &   \\
\tablevspace
 \citet{Unsworth2003} & \textbf{English} German (UK) & 1 & NA & {None}\\
\tablevspace
 \citet{Austin2009} & \textbf{Basque} \textbf{Spanish} (Spanish Basque country) & 20 & – &{None}\\
\lspbottomrule
\end{tabularx}
\caption{Studies investigating interaction in \textbf{Root} \textbf{Infinitives} in single-language utterances}
\label{tab:vihmannvihmann:7}
\end{table}


\textit{Root infinitives} (RIs, \tabref{tab:vihmannvihmann:7}) involve the delayed acquisition of finite verb marking.\footnote{We distinguish ‘inflection’ from ‘root infinitives’ based on the terms used in the different studies.} They are the main exception to the general finding of interaction: Here, none of the studies found any such evidence. \citet{HulkMüller2000} argue that although Root infinitives satisfy one of the two conditions they posit for cross-linguistic influence (involving the interface between syntax and pragmatics), they fail to fulfill the second condition, with cross-linguistic ‘overlap’ such that ``the input of one of the adult languages reinforces a misanalysis of RIs (in declarative root clauses) as correct in the other language" (p. 240). In a close study of one Dutch-French and one German-Italian bilingual child, Hulk and Müller find that the rate of use of RIs falls within the range of monolingual children for each of the languages, although they are rare in Italian.


\citet{Unsworth2003} extends \citeauthor{HulkMüller2000}’s analysis to an English-German bilingual child, where a cross-linguistic effect might be expected, given the surface similarity between the two languages: Bare verb forms occur in German input as well as in English. Unsworth’s bilingual participant, growing up in Germany, produces too few English bare forms for analysis, but no effect of English on her German is observed with respect to root infinitives: ‘Annie’s German RIs exhibit the same referential properties and contain the same predicate types as her monolingual peers’ (p. 154). Unsworth suggests that the ‘overlap’ condition must mean ‘overlap between forms in the input of the two languages’, and specifically, partial overlap.\footnote{As a reviewer pointed out, in the case of fully overlapping adult structures cross-linguistic effects could not be detected in a child’s production.}  A bilingual child learning German, with its more frequent overtly inflected verbs, could be expected to more easily move beyond the RI stage (i.e., acquire the inflectional markers in English), though Unsworth’s data are insufficient to test the idea. (Compare this with the discussion of the acquisition of finiteness in French and English in \citealt{ParadisGenesee1996}.)

In a comparatively large cross-sectional study, \citet{Austin2009} does find more omission of obligatory verbal inflection (carried by auxiliaries) in the Basque of the younger bilinguals (under age 2;8) than in monolinguals, correlated to some extent with lower MLU, but Austin argues against interpreting this as reflecting influence from Spanish. Both the monolingual and the bilingual children had already advanced beyond the ‘root infinitive stage’ in Spanish, where verbal inflection tends to be learned considerably earlier than the equivalent morphology in Basque. Austin draws instead on \citegen{Gathercole2007} suggestion that a relatively lower amount of input in each of the languages of bilingual children may lead to a temporary delay in some aspects of morphosyntactic acquisition, in comparison with typical monolingual development; a ‘critical mass’ may be needed for the acquisition of some grammatical features. The fact that most of the children showing use of more adult-like morphology in Basque were attending schools where all instruction is provided in Basque (with Spanish taught as a subject) supports this hypothesis; children attending schools where the day is split between instruction in the two languages tended to make use of root infinitives in Basque for a longer period. Parental usage was not monitored. Thus ‘delay’ here, \textit{vis à vis} monolingual expectations, is taken to be the result of the process of bilingual acquisition \textit{per se} (with input split between two languages and therefore reduced in each), rather than being due to any structural characteristic of the competing language.

The studies addressing aspects of \textit{determiner phrases} (DPs, \tabref{tab:vihmannvihmann:8}) are the most heterogeneous in terms of the specific object of analysis.

\begin{table}
\small
\begin{tabularx}{\textwidth}{p{2cm}>{\raggedright}p{2.5cm}rcQ}
\lsptoprule
 {Study} & { Languages}  {(location)} &  {N}  &  {Dominance} &  {Type}  {of}  {Interaction}\\
\midrule
\raggedright\citet{KupischEtAl2002} & \textbf{German} Italian (Germany) & 2 & - & None\\
                         & French \textbf{German} (Germany) & 1 &  & \\
\tablevspace
\raggedright\citet{Kupisch2007} & \textbf{German} Italian (Germany) & 4 & + & Acceleration: In determiners in German \\
\tablevspace
\raggedright\citet{Chang-Smith2010} & \textbf{English} Mandarin (Australia) & 1 & NA & None\\
\tablevspace
\raggedright\citet{EichlerEtAl2013} & French Italian (not specified) & 5 & + & \multirow{4}{3.5cm}{Delay: Errors in neuter marking in German three-way gender system (due to parallel acquisition of Romance two-way system)}\\
                         & French \textbf{German} (Germany) & 5 &  & \\
                         & \textbf{German} Italian (Germany) & 7 &  & \\
                         & \textbf{German} Spanish (Germany) & 2 &  & \\
\tablevspace
\raggedright\citet{LarrañagaGuijarro-Fuentes2013} & \textbf{Basque} \textbf{Spanish} (Spain) & 2 & - & Delay in acquisition of Spanish gender\\
\tablevspace
\raggedright\citet{AnderssenBentzen2013} & English \textbf{Norwegian} (Norway) & 1 & - & Delay in acquisition of Norwegian determiner affix\\
\tablevspace
\raggedright\citet{RodinaWestergaard2013} & English \textbf{Norwegian} (Norway) & 2 & NA & None\\
\tablevspace
\raggedright \citet{HervéSerratrice2018} & \textbf{English} French (England) & 2 & + & Acceleration in Eng. determiner provision Delay in  French plural/ generic determ. provision \\
\lspbottomrule
\end{tabularx}
\caption{Studies investigating cross-linguistic interaction in single-language utterances: \textbf{Determiner} \textbf{phrases}.}
\label{tab:vihmannvihmann:8}
\end{table}

\citet{Kupisch2007}, looking at determiner provision in four children acquiring German and Italian in Germany, found evidence of acceleration in German under the influence of the more accessible determiners in Italian. In contrast, \citet{KupischEtAl2002}, who considered data from both Italian and French children bilingual with German, found later acquisition of gender marking on determiners in the one French child, but ascribed the delay to a lack of sufficient exposure to French. 

\citet{HervéSerratrice2018} found evidence of acceleration in the acquisition of English determiners under the influence of French, even when French was the weaker language. Cross-linguistic transfer was also observed, from English to French, resulting in article omission, but only in periods of English dominance and with considerable lexical specificity.   

\citet{AnderssenBentzen2013} and \citet{RodinaWestergaard2013} analyzed DPs in one and two girls (respectively) learning Norwegian in Norway with English as the home language; they focused on definiteness and gender/declension class, respectively, and found good evidence of interaction in one child, Emma, while in most respects the second child in Rodina and Westergaard’s study resembled the two monolingual children in the study more than she did her younger sister Emma. (Note that in Norwegian definite articles are expressed by inflection rather than by separate determiners, as in the other languages involved here.) 

On the other hand, \citet{EichlerEtAl2013} failed to find evidence of interaction in gender marking in children learning German (with its three genders) and a Romance language (with two) but did observe some delay in the children’s weaker languages. 

The bilingual data in \citet{LarrañagaGuijarro-Fuentes2013} derive from two children learning Spanish alongside Basque, a language that lacks both gender features and nominal agreement. The bilinguals are found to master gender assignment more slowly than the monolingual Spanish-learning child who serves as a point of comparison. The authors tentatively ascribe the delay to Basque influence on Spanish, but they consider their data too limited for a solid conclusion. Finally, \citet{Chang-Smith2010} found no evidence of interaction in the production of classifiers in the single child she observed, acquiring Mandarin alongside English in the United States.

Overall, we find that cross-linguistic interaction is widely reported in the literature on early bilingualism. However, it is not found in all cases, and some authors note that the findings based on small samples are inconclusive. The results are complex, as studies approach the issue from different angles and the evidence may be difficult to evaluate in the absence of baseline or comparative data.

\subsection{Evidence from code-switching or ``mixed-language utterances''} \label{sec:vihmannvihmann:3.2}

We initially identified 21 papers investigating \textit{mixed utterances} in bilingual children. However, many of these were not directed at the question of linguistic interaction and failed to specify the morphosyntactic structures involved in any mixing. Those papers were concerned with other issues, such as evidence of children’s ability to separate their languages or to follow adult-like constraints, the children’s language choices, or the extent to which the input may be the source of children’s mixing. 

Here we restrict ourselves to the ten studies that include the question discussed above: Which language structures are involved in cross-linguistic interaction in the young bilingual child? Although some of these studies specifically pick out aspects of determiner phrases (DPs) for analysis, most fail to focus on any one structure but instead provide analyses and examples of children’s utterances that do or do not follow the constraints found to apply to much of adult code-switching, most often referencing \citegen{Myers-Scotton1993duelling} Matrix Language (ML) framework. However, each of the studies can be broadly categorized as primarily addressing one of three aspects of morphosyntax (see Tables 9–11). Interaction between the children’s two languages is necessarily implicated in these studies, with their attention to mixed-language utterances.\footnote{Under ‘mixed language utterances’ we include only intraclausal or intralexical switches, in which the grammars of the two languages confront each other (\citealt{Myers-Scotton1993duelling}).}

The studies are listed chronologically within each table. The \textbf{community} \textbf{language(s)} are marked in bold (with study location). We indicate whether the authors see Dominance as a factor affecting the interaction in question (+ / -). An important factor that differs across the studies is the question of compatibility between the two languages with respect to the structures in question; this is indicated in the rightmost column, under Compatible structures, showing whether the structures overlap in the languages in question (+ / -).

As regards the characterization of dominance, five studies (\citealt{CantoneMüller2008,EichlerEtAl2012,GaskinsEtAl2021,ParadisEtAl2000,YipMatthews2016}) formally compare the children’s MLU in the two languages (although \citealt{YipMatthews2000} caution that establishing MLU for Cantonese can be problematic); \citeauthor{ParadisEtAl2000} also consider vocabulary and parental report, while \citeauthor{CantoneMüller2008} additionally include Upper Bound, lexical diversity (numbers of noun and verb types) and number of utterances in half-hour recordings. \citet{Soriente2014} and \citet{Vihman2016,Vihman2018} draw on informal observation of the child’s usage and extent of exposure to the two languages. \citet{Gawlitzek-MaiwaldTracy1996} report dominance based on the children’s language preferences (proportion of utterances produced in each language). \citet{BernardiniSchlyter2004} recruited unbalanced bilingual participants to test their ideas; they base their assessment of dominance on extent of input in each language.

\subsubsection{Compatibility of structure} 

Studies investigating cross-language interaction in \textit{single-language} utterances must necessarily select a linguistic phenomenon which differs across the two languages to identify the interaction. Unlike the studies discussed in \sectref{sec:vihmannvihmann:3.1}, all the studies included in \sectref{sec:vihmannvihmann:3.2} do by definition provide evidence of some interaction between languages, given the language mixing in the utterances. An additional question that arises here is the effect, on mixing, of the compatibility or equivalence of the relevant structures in the two languages. Note that it has long been claimed that code-switching requires some degree of compatibility (\citealt{Poplack1980}, \citealt{Myers-Scotton1993duelling}). As noted above, an early suggestion for predicting where interaction could occur in bilingual development was (partial) overlap between the structures of the two languages (\citealt{HulkMüller2000}); for mixing in adult studies, the point was encapsulated in the ‘structural equivalence’ constraint (\citealt{Poplack1980}). Accordingly, we mark the language pairs for compatibility in Tables 9–11. Here we see that about half of the studies investigate intraclausal mixing in structures which are not compatible across the two languages.

We define compatibility as overlap, or similarity, both in linear order and in the type of linguistic elements required for the construction in question. Compatibility is more nuanced than the binary coding suggests, however: For instance, an inflectional paradigm can include areas of overlap (e.g. suffixes are required for third person, present-tense verbs in both English and German) and of non-overlap (first and second person suffixes are required in German but not English; two languages may both require prenominal articles, but categorize them differently for gender, as with French and German). In this section, we have nevertheless marked the structure under investigation as either compatible or not, as this informs both the research question and the potential for structural interaction. Compatibility again implies partial rather than complete overlap (see fn. 5).

\subsubsection{Determiner Phrases}

\tabref{tab:vihmannvihmann:9} gives an overview of the studies in our sample investigating cross-linguistic interaction in determiner phrases in mixed-language utterances. German and Italian share the basic structure Determiner + Noun, with the gender of the noun indexed on the determiner in both languages. \citet{CantoneMüller2008} find that in 5\% of the mixed-language DPs the children encode on the determiner not the gender of the switched noun they produce, but that of the \textit{equivalent noun} in the other language. In other words, the children sometimes access the intended word form in one language but index, on the obligatory determiner, the gender of the equivalent noun in the other language.

\begin{table}
\begin{tabularx}{\textwidth}{QQrcp{1.8cm}}
\lsptoprule
{ Study} & { Languages} { (location)} & { N} & {{ Dominance}} & { Compatible structures}\\
\midrule
\citet{CantoneMüller2008} & \textbf{German} Italian (Germany) & 4 & - & { +}\\
\tablevspace
\citet{EichlerEtAl2012} & French Italian (not specified) & 1 & { -} & +\\
& French \textbf{German} (Germany) & 5 &  &   \\
%\hhline%%replace by cmidrule{~--~~~}
& \textbf{Germa}n Italian (Germany) & 7 &  &   \\
%\hhline%%replace by cmidrule{~--~~~}
& \textbf{Germa}n Spanish (Germany) & 2 &  &  \\
\tablevspace
\citet{GaskinsEtAl2021} & \textbf{French} Russian (France) & 1 & { -}& -\\
& English \textbf{German} (Germany) & 1 &  &   +\\
%\hhline%%replace by cmidrule{~--~~-}
& \textbf{English} Polish (England) & 2 &  & -\\
%\hhline%%replace by cmidrule{----~-}
\lspbottomrule
\end{tabularx}
\caption{Studies investigating cross-linguistic interaction in \textbf{determiner} \textbf{phrases} in mixed language utterances. The community language(s) are marked in bold. We indicate whether the authors see Dominance as a factor affecting the interaction (+ / -) and whether the structures under investigation are seen as partially compatible  (+ / -).}
\label{tab:vihmannvihmann:9}
\end{table}

We give one example with a German noun in \REF{ex:vihmannvihmann:1}, from a balanced bilingual child aged 3;8\footnote{They do not state how they drew the line between the ‘balanced’ and ‘unbalanced’ bilingual children whose data they present.}:

\ea%1
    \label{ex:vihmannvihmann:1}
\gll            {al} \textbf{{schlange}}\footnotemark\\
  to.\textsc{def}.\textsc{m}  snake.\textsc{f}\\%
  \z
\footnotetext{Cf., Italian:
    \ea
    \gll al        serpente\\
    to.\textsc{def.m} snake.\textsc{m}\\
    \glt `to the snake'
        \z
}


Here, then, the DP structures are compatible, but the lexical categorization of gender is not. These errors are more often seen in the productions of children with less balanced bilingualism, but they occur in the stronger as well as the weaker language. While 5\% is a small minority of mixed-language DPs, it is unlikely to come from random errors in gender assignment, as the data rarely show errors with nouns which have the same gender in the two languages. 

\citet{EichlerEtAl2012} also report this minority pattern in their study of 15 children in four different language pairs, with the gender matching the equivalent noun in the child’s other language some 20-25\% of the time, considerably more often than in \citet{CantoneMüller2008}. \citet{EichlerEtAl2012} find that the determiner is always in the language of the verb, which is taken to establish the ML for the clause\footnote{Identification of the ML is a major problem for the analysis, especially in two-word combinations. \citet{EichlerEtAl2012} find that internal evidence – i.e., the language of the verb – is a more reliable anchor for analysis than the language of the discourse context, typically set by the interlocutor but not necessarily respected by the child. (See detailed discussion in \citealt{ParadisEtAl2000}: 251f.; \citealt{BernardiniSchlyter2004} also address the issue.)}; only the noun is switched, and the gender is generally that of the switched noun, not that of the verb and determiner, despite its being marked on the determiner. Balanced and unbalanced bilingual children do not differ in this regard, nor does it matter whether the noun comes from the child’s stronger or weaker language (see \citealt{EichlerEtAl2012}, Figure~6: 250).

Finally, \citet{GaskinsEtAl2021} include three language pairs, only one of which involves compatible DP structures: English and German both require articles before the noun in most DPs, whereas Polish and Russian lack articles. These authors find that, although the two children learning Polish and English do not switch elements within the DP, equivalence between structures cannot constitute an absolute constraint, given that the child learning Russian and French does insert both ‘proto-determiners’ and indefinite articles from French before Russian nouns. The study does not report the language of the clause or utterance in which they occur. As the child advances in her knowledge of Russian, the switches disappear. \citet{GaskinsEtAl2021} suggest that mastering a language with a greater variety of article types (due to case and gender marking) supports earlier ‘segmentation’, or the distinction between determiner and noun within the DP; this distinction facilitates switching. 

\subsubsection{Inflection}

In \tabref{tab:vihmannvihmann:10} we see the extent of cross-linguistic interaction in inflection in mixed-language utterances in the studies in our sample.

\begin{table}
\begin{tabularx}{\textwidth}{QQrcp{1.8cm}}
\lsptoprule
{ Study} & { Languages} { (location)} & { N} & {{ Dominance}} & { Compatible structures}\\
\midrule
\citet{ParadisEtAl2000} & English \textbf{French} (Montreal, Canada) & 15 & - & { +}\\
\tablevspace
 \citet{BernardiniSchlyter2004} & Italian \textbf{Swedish} (Sweden) & 1 & { +}  & -\\
%\hhline%%replace by cmidrule{~--~~~}
& French \textbf{Swedish} (Sweden) & 2 &  &   \\
\tablevspace
 \citet{Soriente2014} & \textbf{Indonesian} Italian (Indonesia) & 1 & { +} & -\\
\tablevspace
 \citet{Vihman2016} & English \textbf{Estonian} (Estonia) & 2 & { -}
& +\\
\tablevspace
 \citet{Vihman2018} & English \textbf{Estonian} (Estonia) & 2 & { -}
& -\\
\lspbottomrule
\end{tabularx}
\caption{Studies investigating cross-linguistic interaction in \textbf{inflection} in mixed language utterances}
\label{tab:vihmannvihmann:10}
\end{table}

In a study of 15 French-English bilingual children living in the francophone community of Montreal, \citeauthor{ParadisEtAl2000} (\citeyear{ParadisEtAl2000}; see also \citealt{ParadisGenesee1996}) set out to determine the extent to which the children’s code-switching adheres to the structural constraints widely claimed to govern adult switches (\citealt{Myers-Scotton1993duelling}). Analyses of the children’s mixed-language utterances in recordings made at ages 2;0, 2;6, 3;0 and 3;6 also made it possible to test \citegen{Meisel1994} proposal that in the earliest period of word combination linguistic principles would not yet apply, but that once children have developed productive use of functional categories, both single- and mixed-language utterances could be expected to conform to linguistic principles. A developmental difference in the acquisition of verbal markers of tense and agreement between French (early) and English (later) provides a useful testing ground for these ideas.

Use of French system morphemes (e.g., auxiliary verbs, copulas, modals, tense and agreement inflections) in English clauses proved to be relatively frequent (at 20-30\%) in the second and third recordings; the more balanced bilinguals produced proportionately more such switches (once their utterances had become more advanced or ambitious), in accord with the ‘bootstrapping’ hypothesis of \citet{Gawiltzek-MaiwaldTracy1996}. The authors find no evidence of a sharp developmental shift, such as Meisel predicted; \citeauthor{Myers-Scotton1993duelling}’s system-morpheme constraint was more often observed (82\%) than violated across all four sessions for the 15 children. (Note, however, that \citeauthor{ParadisEtAl2000} determined ML based on the developmental psycholinguistic criterion, excluding mixed utterances from the calculation, rather than on a clause-by-clause basis, as has been done in more recent work.) \citeauthor{ParadisEtAl2000} conclude that ‘it is the children’s lexicons that must develop in order to give them the tools to adhere more strictly to the constraint’ (\citeyear{ParadisEtAl2000}: 259).

The ‘ivy hypothesis’ (\citealt{BernardiniSchlyter2004}) further elaborates this line of thinking. Focusing on three children with more advanced development in Swedish than in their other language, French or Italian, \citet{BernardiniSchlyter2004} find that most switching involves clausal structure in Swedish with ‘lower’ elements taken from the less dominant, less developed language. The authors conclude that ‘practically all mixed utterances produced by these children during the period studied were such that the “missing” elements – which represent portions of higher syntactic structure – were “replaced by” their counterparts in the Stronger Language’ (p. 65f.). This is the case despite the fact that Swedish, as a V2, non-pro-drop language, has a clausal structure that is incompatible with those of the Romance languages.

The case study \citet{Soriente2014} presents involves a still greater imbalance between the child’s two languages, Jakarta Indonesian, the language of the community where he is being brought up, and Italian, his mother’s language. The two languages are structurally very different: Italian requires suffixal markers of tense and aspect whereas these concepts are marked lexically, and only optionally, in Indonesian, where the pragmatics of the situation often replace overt marking. Indonesian was the dominant language, used by a full-time caregiver as well as the father. The child was found to first incorporate Indonesian aspectual function words into his Italian, and then to produce Italian word forms that were the rough equivalent of those words, before he began to produce actual inflections in Italian. For example, the earliest-learned Indonesian verbal aspect marker, perfective \textit{(s)udah,} was frequently inserted into Italian clauses, as shown in \REF{ex:vihmannvihmann:2a}; later, Italian \textit{basta} ‘enough’ was used to express past tense, taking the place of \textit{udah}, as in \REF{ex:vihmannvihmann:2b}.

\ea%2
    \label{ex:vihmannvihmann:2}
\ea \label{ex:vihmannvihmann:2a}
\gll \textbf{{udah}} {volato} (2;9)\\
\textsc{perf}     flown\\
\glt `[It] flew away'

\ex \label{ex:vihmannvihmann:2b}
\gll   {Gulli} \textbf{{basta}} fare   il bagno   con   Papà. \\
 G.  enough  do.\textsc{inf}  the bath  with  Daddy  \\
 \glt `Gulli \textbf{enough} had a bath with Dad'
    \z
\z

Given the relative transparency of the Indonesian system, in addition to the child’s greater exposure to input in that language, the uses the child makes of the ‘scaffolding’ that Indonesian can supply are understandable and in agreement with the bootstrapping idea. The advance from switching to loan translation or calques before the child finally begins to acquire Italian inflections is an interesting further step, one that has seldom been reported.

\citet{Vihman2016,Vihman2018} looks at diary data in two siblings exposed to both Estonian and English in the home and Estonian in the larger community, including preschool. Despite the probable imbalance in input in favor of Estonian, the children mostly spoke in English to one another; the two languages were largely in balance for both. The data, which derive primarily from the younger child at around age three, provide evidence of intense co-activation between the two languages and interaction of the grammatical systems, with use of both bare noun and verb stems and switches in inflectional marking in both languages. 

\citet{Vihman2016} focuses on verbal morphology. Although the languages are to some extent compatible here, as both mark tense with suffixes, Estonian has complex morphology, often including stem allomorphy, and generally expresses tense and person with separate suffixal morphemes. In general, these diary data draw on the daily life of the child, while also being avowedly biased toward more complex or puzzling instances, or (as \citealt{Serratrice2005} suggests with reference to another diary study) more creative uses. Vihman points up the many types of morphological patterns in the switching that she observes, including verbs (a) marked with inflections from either language, (b) lacking appropriate inflectional morphology altogether, and (c) using ‘blended’ morphemes that draw on points of phonological similarity in the affixes of the two languages.

\citet{Vihman2018} addresses the incompatibility of the two languages in many aspects of nominal morphology, based on switches recorded in the same diary material. Estonian requires extensive case marking, often with stem changes; this creates incompatibilities with switched nouns, with some switched as bare stems and left unmarked, some bearing marking in the language of the verb, others in that of the borrowed noun, and yet others in both. For example, for a period between age 3;5 and 4;1 the younger child frequently produced multiple switches within a single utterance, sometimes producing an Estonian verb and direct object with an intervening English determiner, as in example \REF{ex:vihmannvihmann:3}, where the object is unmarked, though it would require marking in Estonian:

\ea%3
    \label{ex:vihmannvihmann:3}
           {there was a big girl and she}
     \gll \textbf{{veereta-s}} {a} \textbf{{pall}} (3;11.24)\\
        roll-\textsc{3sg.pst}   ball.\textsc{nom.sg}\\
\glt `There was a big girl and she rolled a ball'
\z

Vihman argues that the basic Matrix Language Framework assumption that the two languages involved in code switching are necessarily asymmetrical, with just one of them providing the grammatical frame, cannot be defended here, given the dynamic nature of emergent grammars. She also provides examples of convergence in single-language utterances of the same two girls, with loan translations or calques; the transfers are bidirectional (but not quantified for frequency). This leads her to conclude that, ‘for bilingual children, we cannot assume that their utterances are framed according to a grammar equivalent to that of either target language’ (p. 18). 

\subsubsection{Word order}

\tabref{tab:vihmannvihmann:11} shows the cases of cross-linguistic interaction in word order that the studies in our sample reported for mixed-language utterances.

\begin{table}
\begin{tabularx}{\textwidth}{QQrcp{1.8cm}}
\lsptoprule
{ Study} & { Languages} { (location)} & { N}  & { Dominance} & {{ Compatible structures}}\\
\midrule
\citet{Gawlitzek-MaiwaldTracy1996} & English \textbf{German} (Germany) & 1 & { +} & -\\
\tablevspace
\citet{YipMatthews2016} & \textbf{Cantonese} \textbf{English} (Hong Kong) & 9 & + & { +}\\
\tablevspace
\citet{Vihman2018} & English \textbf{Estonian} (Estonia) & 2 & - & { -}\\
\lspbottomrule
\end{tabularx}
\caption{Studies investigating cross-linguistic interaction in \textbf{word} \textbf{order} in mixed language utterances}
\label{tab:vihmannvihmann:11}
\end{table}

In their case study of a child living in Germany with exposure to English and German, \citet{Gawlitzek-MaiwaldTracy1996} argue that code-switching provides a ‘temporary relief strategy’ for children still in the process of developing their morphosyntactic knowledge; they see language mixing as helping the child to bridge structural as well as lexical gaps (‘syntactic bootstrapping’). They point out that English and German are sufficiently similar to ‘invite the child to assume parallel syntactic formats’ (p. 906), as in \textit{Ich} \textbf{\textit{cover} }\textit{mich}\textbf{\textit{self} \textbf{up} }‘I cover myself up’ (produced at 2;6). Yet mixing also occurs in clauses that require different word orders in the two languages, as in \textbf{\textit{Clean}}\textit{st-du dein} \textbf{\textit{teeth}}? ‘Did you clean your teeth?’ (2;9), where the subject-verb inversion of a main verb observes German but not English clausal structure. The suggestion is that earlier development of constructions in one language allows the child to ‘pool her resources’, ‘taking and combining what is available to her in both languages, in a lexical as well as in a structural sense’ (p. 920). Among her examples of convergence in word order in single-language utterances, \citet{Vihman2018} includes production, in English-only clauses, of the inverted object-verb-subject order that occurs in Estonian object-initial syntax. This similarly suggests that grammatical structure may draw on resources from either language. It casts the bilingual utterances in a new light: We have coded ‘compatibility’ with reference to adult structures, but that need not constrain code-switching in early bilingual productions. 

\citet{YipMatthews2016} address ‘code mixing’ in children learning Cantonese and English in Hong Kong, where switching between the two languages is the community practice. They point out that the congruence in word order between the two languages facilitates switching; this includes the analytic nature of both grammars. For all nine children switching occurs most frequently in the ‘Cantonese context’ (i.e. when the person recording them is using Cantonese), although most of the children also mix in the English context and most show considerable variation across the one- to two-year span of the recordings. In English utterances the addition of a Cantonese sentence-final particle (glossed SFP) is the most common switch, as in \REF{ex:vihmannvihmann:4}:

\ea%4
    \label{ex:vihmannvihmann:4}
\gll      You   tidy   up \textbf{{laa1}}\footnotemark{}\\
you   tidy   up   \textsc{sfp}\\
\glt `You tidy up.' (3;6;8, English context)
\footnotetext{The numbers following Cantonese words mark the lexical tones.}
\z

The children may feel that something is missing in their English utterances in the absence of such completive particles. The authors note that these particles are also switched by adults in Singapore English, for example.

Unlike what has been found in European languages (e.g., German/ Italian: \citealt{Cantone2007}: 173), verbs are the part of speech most often switched, typically from English into Cantonese. In these cases Cantonese aspect markers (ASP) are added to the English verbs, as shown in \REF{ex:vihmannvihmann:5}.

\ea%5
    \label{ex:vihmannvihmann:5}
\gll           {dim2gaai2} \textbf{{turn}} {zo2?}\\
why              turn   \textsc{asp}\\
\glt `Why did it turn?' (3;4;18, Cantonese context)
    \z

Bare stems in Cantonese are also mixed into English utterances, but the addition of English tense-aspect morphology is rare.

Another frequent type of switch, shown in \REF{ex:vihmannvihmann:6}, involves English verb particles, sometimes with Cantonese material intervening between the verb and its particle:

\ea%6
    \label{ex:vihmannvihmann:6}
\ea
\gll   Dim2   gaai2   lei5 \textbf{{throw}} ni1   go3 \textbf{{away}} {ge3?}\\
how   come   you   throw   this   CL\footnotemark{}    away \textsc{sfp}\\
\glt  `Why are you throwing this away?' (3;6, Cantonese context)
\footnotetext{CL is presumably ‘classifier’, but Yip and Matthews do not spell it out.}
\ex
\gll Ji1   zek3 \textbf{{slide}} m4   dou2 \textbf{{down}}\\
this   \textsc{cl}  slide   not   succeed   down\\
\glt `(With) these (shoes on) you can’t slide down.' (2;6, CC)
    \z
\z

\citeauthor{YipMatthews2016} see these split-verb switches, unattested in adult code switching, as reflecting ‘cross-linguistic influence, which takes place between the developing grammars, independent of code-mixing’ (p. 12). In other words, above and beyond the fact of producing a mixed-language utterance, they see the insertion of a Cantonese constituent between the English verb and participle as an effect of English grammatical structure. The authors also find evidence of the English split verb-particle structure in the children’s single-language utterances: 

\begin{quote}
The code-mixed cases…are thus congruent with the children’s developmental grammar for Cantonese, if not with the target grammar of Cantonese (p. 12).
\end{quote}

This point is echoed in \citet{Vihman2018}, in relation to the convergence found in the bilingual siblings’ single-language structures.

\section{Discussion and conclusions}

In their meta-analysis of experimental studies with older children than those included in the studies we review here, \citet{VanDijkEtAl2022} find that interaction between the languages of a bilingual child neither diminishes nor increases with age; they conclude that cross-linguistic influence is ‘part and parcel of bilingual development’ (p. 902). In our overview of studies of emergent morphosyntax as observed in children aged about 2 to 4 years, we have similarly found that most studies report evidence of linguistic interaction, and that this affects all but one of the structures investigated. In contrast with the findings of Van Dijk et al., however, many of these studies report a decrease in cross-linguistic effects as the children gain mastery of the two languages. Furthermore, a delay in the acquisition of certain structures, compared to monolinguals, often appears to be due to reduced exposure to one of the languages, or of reduced opportunities for its use – due, in other words, to the bilingual experience itself rather than to any specific structural effects. This may be an effect of bilingual language use that is particular to the earliest period of grammatical development.

\citet{VanDijkEtAl2022} also find that neither surface overlap between syntactic structures nor language domain (such as the interface between syntax and pragmatics) are significant predictors of interaction, contrary to much of the early literature. \citet{Serratrice2013}, in a qualitative review of the literature to that date, also found counterevidence to the hypothesis that only overlapping structures or interfaces can account for the linguistic interaction identified in the speech of bilingual children. She cites evidence and arguments to suggest that the nature of bilingual processing and use may help to better account for the full range of cross-linguistic influence observed in younger children, both in small-N studies and in larger studies that include tests of comprehension. In other words, it is suggested here again that the bilingual experience itself leads to cross-linguistic effects in children, as it does in adults (\citealt{Kroll2013}).

RQ1: \textit{Which linguistic structures are involved in cross-linguistic interaction in the young bilingual child?}

Our first research question concerned the aspects of grammatical systems likely to be affected by cross-linguistic interaction. To address this, we looked at 47 studies of early bilingual production and attempted to establish which structures are involved in either suspected cross-linguistic influence or inclusion of surface elements of both languages in a single utterance (code-switching). Overall, we found that most of the structures we considered were involved in bilingual interaction of one kind or another. Those structures cover most aspects of the grammatical systems acquired in the first years of language use: determiner use in nominal systems, inclusion of copulas in predicates, subject and object provision, inflection in both noun and verb systems, wh-movement and word order. Accordingly, the basic finding is that almost every aspect of the early developing grammar is open to the influence of bilingual interaction.

RQ2a: What \textit{types of interaction} have been reported?

Tables 1-8 identify findings of acceleration (three structures and four studies), delay (six structures, 20 studies) and transfer (two structures, three studies). These results cannot be deemed definitive, given the small numbers of children involved and the various complications affecting analysis of these few cases (issues around the extent of input and practice the children may have experienced in one of the languages, above all, but also the indisputable effect of individual differences in learning rates, strategy, and so on). However, the overview is certainly indicative of the kinds of interaction that may be found between languages that contrast with respect to a range of morphosyntactic structures. We hope this overview of the literature to date will encourage further systematic study, including studies of larger numbers of younger children learning language pairs that differ in selected ways, to supplement the experimental studies of older children.

RQ2b: \textit{And to what extent does dominance in one language account for interaction effects?} 

In \citet{VanDijkEtAl2022}, language dominance proved to be the single most significant predictor of interaction – but only when defined in terms of the \textit{relation of the language tested experimentally to the ‘societal language’}, or the language of the community in which the child is living. That is, there was more evidence of interaction in children tested in a language other than that of the community – the ‘weaker’ language, by that definition – although some effects were bidirectional. When dominance was instead defined and measured in ways particular to each study (for these older children the measures included amount of exposure or use, size of lexicon and fluency ratings by caretakers), dominance was not a significant predictor. Van Dijk et al. interpret their failure to find a significant effect in this analysis as likely due to differences in the ways the authors of various studies arrived at their categorization. In the case of children under age four, who may or may not have been regularly exposed to a range of speakers for each of their two languages, the societal language is likely to be less directly connected to relative knowledge or proficiency.

In the literature on younger children, where the issue has long been a subject of debate, dominance is generally measured by Mean Length of Utterance (MLU) in each language, although numerous other measures are used in the studies examined here as well, as discussed in 3.1 and 3.2. Of the 29 studies that mention it, 13 see dominance as a factor. Note that in their account of gender assignment in determiners, \citet{EichlerEtAl2013} failed to observe cross-linguistic effects but saw dominance as interacting with the overall ranking of languages in terms of children's accuracy in gender assignment. 

On the other hand, at least three studies (\citealt{AndersenBentzen2018,HulkMüller2000,Muller2001}) specifically argue that only internal structural issues are at play; they judge that dominance cannot be a key factor, since they find interaction effects to be bidirectional (cf. also \citealt{AdnyaniEtAl2018}, who also find bidirectional effects, but with more instances of influence from the dominant to the non-dominant language than vice versa). Given the differences in definition and measurement of dominance across the studies and language pairs we have considered, we can draw no solid conclusions. Ultimately, the effect of dominance may depend more on individual differences, and the shifting strength of a child’s proficiency in each language, than on any one bilingual combination or situation.

It is worth noting that the linguistic experiences of children sharing a language pair, even in the same community, may be quite different, depending not only on differences in family language policy and usage but also on the individual child’s rate of language development and overall maturation. In general, these studies point up the importance of individual differences, particularly in the early years, when children acquire morphosyntactic knowledge at very uneven rates; comparisons of longitudinal studies of one or two children make this particularly evident (\citealt{Serratrice2005}). At older ages, developmental differences begin to even out and the societal (or school) language typically becomes the stronger one.

A nice example of the difficulty of drawing broad conclusions from studies of one or two children is provided by \citet{AndersenBentzen2018}, two of whose participants were siblings, 10 years apart in age (but followed over the same early developmental period; cf. also \citealt{RodinaWestergaard2013}): The older child showed little cross-linguistic interaction, whereas the younger child, said to be well balanced in her two languages, showed considerable interaction in both directions. For example, Emma showed the influence of Norwegian in her question formation, as in example \REF{ex:vihmannvihmann:7}:

\ea%7
    \label{ex:vihmannvihmann:7}
     \textit{Drive daddy me to} \textbf{\textit{barnehage}}? \\
\glt `Will/did daddy drive me to nursery?'

(\citealt{AndersenBentzen2018}: 8)
    \z

Emma seemed to find language learning more challenging than her older sister had; she may also have had more ‘tolerance for variability’ (\citealt{CattsDavis1984}), or in other words, more willingness to produce target-like and non-target-like structures in parallel.

RQ3: \textit{What role is played by the relative accessibility and compatibility of structures?}

Few of the studies reviewed here mention the possible role of the relative accessibility of structures in the children’s two languages. Two studies – \citet{Kupisch2007} and \citet{HervéSerratrice2018} – report accelerated learning under the influence of the child’s other language, but in an earlier study \citet{KupischEtAl2002} found no interaction between German and a Romance language with respect to acquisition of the same structure (determiner phrases; cf. also \citealt{EichlerEtAl2012}). Both \citet{KupischEtAl2002} and \citet{Austin2009} find that a lack of sufficient exposure to one of the languages is the more plausible explanation for the effects they identify.

Recall that among the conditions that \citeauthor{HulkMüller2000} proposed as being necessary for syntactic cross-linguistic influence to occur was ‘a certain overlap of the two systems at the surface level’ (\citeyear{HulkMüller2000}: 228). This requirement has since been somewhat reformulated but not overturned. The overlap refers to (partial) similarity in input structures in the child’s two languages (\citealt{Unsworth2003}). More specifically, \citeauthor{HulkMüller2000} proposed that an ambiguous structure that allows for more than one syntactic interpretation in language A may correspond to a structure with only one analysis in language B; cross-linguistic influence occurs when the child takes the analysis in language B as evidence for language A. The studies we review investigate structures with various types of (partial) overlap, including similar elements with different word order (e.g. with wh-movement), similar word order but variation in which linguistic categories are expressed (e.g. gender on determiners), similar inflectional categories but differences in distribution (e.g. verb inflection), and structures with optional expression in one language but not the other (e.g. argument omission).

The structures in which cross-linguistic influence has been documented suggest that it is not limited to those with differences in the analyses available in each language, as \citeauthor{HulkMüller2000}’s formulation would have it. From Tables 9-11 we see that many structures that lack structural compatibility nevertheless appear in bilingual children’s mixed-language utterances. \citegen{Soriente2014} example of the development of aspect marking, lexical in Indonesian and inflectional in Italian, provides a good case in point: The child made use of Indonesian lexical aspect to scaffold the development of aspectual marking in Italian. There is ‘overlap’ only in the fact of overt expression of aspect in both languages. The overrepresentation, in the literature, of European languages (\citealt{Kidd2022}; \citealt{YipMatthews2022}), which tend to be both genetically and typologically similar (see \citealt{Haspelmath2001}), may have led to a bias in the evidence with regards to the limits of cross-linguistic influence, giving greater weight to mutual compatibility of structures. To more fully test the effects of structural differences in the input we need more studies of both monolingual and bilingual development, in a wider range of languages representing greater typological diversity.

As the abstract categories of the target grammars cannot be said to be fully available to children (e.g., \citealt{Ambridge2020}), it is unclear how similarity of structure may operate in children’s bilingual production and developing grammars. Experimental and corpus studies of syntactic priming with adults (e.g. \citealt{Kootstra2012,Travis2017}) have shown that representations of syntactic structure may be mapped across languages, but that this requires similarity of both constituent order and syntactic structure. Whereas linear order is undeniably available to children, it is unclear to what extent the more abstract syntactic similarity of constructions may influence their usage, or at what age (for syntactic priming in monolingual children see, for example, \citealt{Rowland2012,PeterEtAl2015,BraniganMcLean2016}). Some studies have found cross-language effects of structural priming in children (e.g. \citealt{Herve2016,Vasilyeva2010,Wolleb2018}), but more work is needed to disentangle the interaction of cognitive and linguistic development with cross-linguistic influence.

Children’s developing grammars are a moving target: Children are continuously (i) reshaping their grammatical knowledge, (ii) advancing in their ability to abstract structure out of learned linguistic material, and (iii) developing their ability to produce structures accordingly. Hence, it may be impossible to define what kind of overlap or how much of it there should be between languages to trigger interaction; the overlap itself, as represented in the child’s mind, may undergo continual shifting.

\section{Future directions}

The attempt to define children’s productions, especially utterances with code-switching, according to target language may be misguided. Children’s grammars are in flux, both developmentally and through ongoing interaction with the other language. Hence, it is not always appropriate (or even possible) to define a matrix language according to the adult grammar (\citealt{Vihman2018}) or to assume a target structure for children’s productions (\citealt{YipMatthews2016}). Nevertheless, the decades of fruitful research included in our review have resulted in a sizable body of evidence that language interaction does take place, in a wide range of structures and language pairs.

The approach of translanguaging studies (e.g. \citealt{Wei2018,Otheguy2015,Balam2021}) is congenial to the view that children make use of their full linguistic repertoires, often without drawing clear lines between their named languages. The translanguaging literature, drawing on a wider array of bilingual speakers – especially older children in educational contexts – provides a useful reminder that bilinguals’ linguistic resources may amount to more than the sum of their parts. That is, use of their languages in combination may be taken to increase expressive capacity. The ‘moving target’ of emerging grammars mentioned above can be seen to continue to develop as children enter school, while the relative dominance of their languages, as well as the repertoire available to them, may shift. We hope that future research will investigate that trajectory, beginning with the early years, when children have yet to gain a full grasp of either grammatical system, through to a context where they have become proficient and educated in one language and may experience a range of attitudes toward the other, in and out of school and in the home.

Bearing in mind the dynamics of emergent grammatical structures and the shifts in the extent or type of exposure to and opportunities for use of each language over the course of development, we cannot with any confidence pin down the developmental points most likely to reflect interaction between the child’s languages, intriguing though that issue might be. Recall the conclusion of \citet{ParadisEtAl2000}, that it is lexical advance that paces the development of morphosyntactic structure – a conclusion that aligns well with the long-standing views and findings of Elizabeth Bates and her colleagues with respect to monolingual acquisition (cf. \citealt{Bates1999}).

Both generative and usage-based accounts have grappled with Root infinitives in monolingual children’s language development because the data are hard to reconcile with the theoretical predictions. It is perhaps unsurprising that this is the only linguistic domain in which none of the studies we reviewed report cross-linguistic interaction. As various factors may be at play in the RI conundrum, an approach that looks closely at the structural differences in the languages, as well as the frequency of occurrence in the input of the syntactic structures in question, the word forms themselves and collocational information, would considerably advance our understanding of optional infinitive usage in bilingual development. \citet{FreudenthalEtAl2024} provide a promising learning model that could readily be tested with bilingual children.

Another factor worth investigating in future research is the question of patterns of use in the wider community (see \citealt{BalamEtAl2021} and \citealt{vanOsch2023}). The extent to which the level and form of bilingualism in the community affects a child’s usage is underexplored in general, but particularly in studies of younger children. The effects of community use may differ by child, by age and by caregiving context. Attitudes toward each language and toward language mixing, the broader prestige of the languages, and the extent to which they are used in similar contexts by the same interlocutors may all play a role in children’s acquisition of the systems.

A fundamental unresolved issue underlying any aspect of cross-linguistic interaction, whether in single- or mixed-language utterances, is that of children’s memory of, or access to, the speech patterns they have heard and practiced, or in other words, the nature of their linguistic representations. We must assume, based on the evidence we have reviewed here, that those representations are variable, not only over developmental time but also over moments of utterance; in addition, we can assume that children experience fluctuations in the relative control available for shaping their utterances to fit these representations, that is, in their control over the process of speech production. Theoretical models developed to account for the intricacies of adult grammar need to incorporate evidence from both language acquisition and bilingual processing. As the increasing pace of studies in this area suggests, the bilingual child can provide an intriguing window into the ‘inner life’ of language in development. 

\sloppy\printbibliography[heading=subbibliography,notkeyword=this]
\end{document}
